\chapter{Dasar Teori}

% ==========================================
% CLASSIFICATION
% ==========================================
\section{Classification}
Klasifikasi merupakan salah satu teknik \textit{supervised learning} yang bertujuan untuk memprediksi label atau kategori dari sebuah sampel berdasarkan fitur yang dimiliki. Model klasifikasi belajar dari data berlabel, kemudian menghasilkan suatu fungsi yang mampu menentukan kelas dari data baru yang belum pernah dilihat [\ref{modulPraktikum}]. 

Sebelum tahapan klasifikasi dilakukan, kualitas data mentah perlu dievaluasi terlebih dahulu melalui proses \textit{Exploratory Data Analysis} (EDA) untuk menemukan pola, mendeteksi anomali, dan menguji asumsi menggunakan statistik deksriptif maupun visualisasi grafis [\ref{tukey}]. Setelahnya, data harus melalui tahap \textit{data preprocessing} dengan memanfaatkan ekosistem pustaka komputasi Python (seperti Pandas dan NumPy) guna membersihkan serta mentransformasikan struktur data agar siap diproses oleh model komputasi [\ref{mckinney}].

Secara umum, klasifikasi kemudian bekerja dengan cara mengekstraksi dan mempelajari pola dari data historis yang sudah bersih tersebut untuk menentukan keputusan kelas yang paling mendekati. Teknik klasifikasi banyak digunakan dalam berbagai aplikasi dunia nyata, seperti deteksi spam pada email, diagnosis penyakit dalam bidang medis, hingga sistem rekomendasi konten.

% ==========================================
% ALGORITMA CLASSIFICATION
% ==========================================
\section{Algoritma Classification}
Terdapat berbagai macam metode yang dapat digunakan untuk melakukan klasifikasi. Berikut adalah beberapa algoritma populer yang sering diimplementasikan dalam pustaka \textit{machine learning} (seperti Scikit-Learn) karena kesederhanaan, efektivitas, dan kecocokannya untuk berbagai jenis dataset [\ref{modulPraktikum}, \ref{geron}]:

\begin{enumerate}
    \item \textbf{Logistic Regression}\\
    \textit{Logistic Regression} merupakan algoritma klasifikasi berbasis regresi linier yang memodelkan probabilitas suatu sampel termasuk ke dalam kelas tertentu menggunakan fungsi sigmoid. Output dari model ini berada dalam rentang 0 hingga 1, yang kemudian dikonversi menjadi label kelas menggunakan ambang batas (\textit{threshold}). Meskipun sederhana, algoritma ini sangat efisien dan bekerja dengan baik pada data yang memiliki hubungan \textit{linear separable}.

    \item \textbf{K-Nearest Neighbors (KNN)}\\
    \textit{K-Nearest Neighbors} mengklasifikasikan suatu data baru berdasarkan mayoritas kelas dari $k$ tetangga terdekatnya. Kedekatan antar data dihitung menggunakan metrik jarak tertentu, seperti \textit{Euclidean distance}. Algoritma KNN bersifat \textit{non-parametric} dan sangat bergantung pada struktur distribusi data serta nilai hiperparameter $k$ yang dipilih.

    \item \textbf{Decision Tree}\\
    \textit{Decision Tree} bekerja dengan cara membagi dataset secara rekursif berdasarkan fitur yang memberikan informasi terbaik. Proses pemilihan fitur pada setiap percabangan (\textit{node}) dilakukan menggunakan ukuran kemurnian seperti \textit{Gini impurity} atau \textit{entropy/information gain}. Kelebihan utama metode ini adalah interpretasinya yang sangat mudah dipahami layaknya aturan logika bersyarat, tetapi mudah mengalami \textit{overfitting} tanpa pengaturan parameter yang tepat (\textit{pruning}).

    \item \textbf{Random Forest}\\
    \textit{Random Forest} adalah algoritma \textit{ensemble} yang menggabungkan banyak \textit{Decision Tree} melalui teknik \textit{bagging} (Bootstrap Aggregating). Setiap pohon keputusan dilatih pada subset data dan subset fitur yang berbeda secara acak, dan hasil prediksi akhirnya ditentukan melalui sistem pemungutan suara terbanyak (\textit{majority voting}). Metode ini membuat model lebih stabil, secara signifikan mengurangi risiko \textit{overfitting}, dan meningkatkan akurasi dibandingkan satu pohon keputusan tunggal.

    \item \textbf{AdaBoost (Adaptive Boosting)}\\
    \textit{AdaBoost} bekerja dengan menggabungkan beberapa \textit{weak learner} (biasanya berupa \textit{Decision Tree} kecil atau \textit{stump}) secara berurutan. Pada setiap iterasi, bobot sampel data diperbarui sehingga model berikutnya akan lebih fokus memperbaiki prediksi pada data yang sebelumnya salah diklasifikasikan. Dengan pendekatan adapatif ini, AdaBoost mampu meningkatkan performa model secara signifikan, terutama pada dataset yang memiliki kompleksitas tinggi.
\end{enumerate}

% ==========================================
% METRIK EVALUASI CLASSIFICATION
% ==========================================
\section{Metrik Evaluasi Classification}
Dalam \textit{machine learning}, khususnya pada pemodelan klasifikasi, diperlukan metrik evaluasi yang terukur untuk menilai seberapa baik model mampu memprediksi label dengan benar. Evaluasi ini dilakukan dengan membandingkan hasil prediksi model terhadap nilai aktual (kebenaran lapangan) melalui sebuah struktur yang disebut \textit{confusion matrix} [\ref{modulPraktikum}, \ref{bishop}].

\subsection{Confusion Matrix}
\textit{Confusion matrix} terdiri dari empat komponen penting yang merepresentasikan persilangan antara prediksi dan nilai aktual, antara lain:
\begin{itemize}
    \item \textbf{True Positive (TP)}: Prediksi kelas positif yang benar secara aktual.
    \item \textbf{True Negative (TN)}: Prediksi kelas negatif yang benar secara aktual.
    \item \textbf{False Positive (FP)}: Prediksi kelas positif yang salah (aktualnya negatif).
    \item \textbf{False Negative (FN)}: Prediksi kelas negatif yang salah (aktualnya positif).
\end{itemize}
Berdasarkan nilai-nilai tersebut, berbagai metrik evaluasi dapat dihitung untuk memahami performa model secara lebih komprehensif.

\begin{figure}[H]
    \centering
    \includegraphics[width=0.6\textwidth]{21.png}
    \caption{Visualisasi Struktur Confusion Matrix}
    \label{fig:confusion_matrix_theory}
\end{figure}

\subsection{Accuracy, Precision, Recall, dan F1-Score}
\begin{enumerate}
    \item \textbf{Accuracy}\\
    \textit{Accuracy} merupakan metrik paling dasar yang menunjukkan proporsi prediksi yang benar dari seluruh total prediksi yang dibuat oleh model. Rumusnya didefinisikan sebagai berikut:
    \begin{equation}
        Accuracy = \frac{TP + TN}{TP + TN + FP + FN}
    \end{equation}

    \item \textbf{Precision}\\
    \textit{Precision} mengukur tingkat ketepatan model dalam melakukan prediksi positif, yaitu rasio berapa banyak prediksi positif yang benar-benar positif secara aktual dibandingkan total prediksi positif yang dibuat. Rumusnya adalah:
    \begin{equation}
        Precision = \frac{TP}{TP + FP}
    \end{equation}

    \item \textbf{Recall}\\
    \textit{Recall} (atau Sensitivitas) mengukur kemampuan model dalam menemukan seluruh sampel yang benar-benar positif. Ini merupakan rasio prediksi positif yang benar dibandingkan total data aktual yang bernilai positif:
    \begin{equation}
        Recall = \frac{TP}{TP + FN}
    \end{equation}

    \item \textbf{F1-Score}\\
    Karena \textit{precision} dan \textit{recall} seringkali bertolak belakang dan memiliki fokus yang berbeda, \textit{F1-Score} digunakan sebagai metrik yang menyeimbangkan keduanya. Metrik ini sangat berguna ketika dataset bersifat \textit{imbalanced} (tidak seimbang) dan diperlukan keseimbangan antara kemampuan model dalam menghindari \textit{false positives} serta \textit{false negatives}. F1-score merupakan \textit{harmonic mean} dari precision dan recall, dengan rumus:
    \begin{equation}
        F1\text{-}Score = 2 \times \frac{Precision \times Recall}{Precision + Recall}
    \end{equation}
\end{enumerate}

\chapter{Hasil dan Pembahasan}

\section{Hasil Praktikum}

\subsection{Impor Pustaka}
\begin{figure}[H]
    \centering
    \includegraphics[width=0.6\textwidth]{1.png}
    \caption{Source Code Impor Pustaka}
    \label{fig:1}
\end{figure}
Penjelasan Gambar :
\begin{itemize}
    \item \textbf{Input:} Mengimpor pustaka inti \texttt{pandas}, \texttt{numpy}, \texttt{matplotlib.pyplot}, serta sejumlah modul dari \texttt{sklearn} (\texttt{train\_test\_split}, \texttt{LabelEncoder}, \texttt{StandardScaler}, metrik evaluasi seperti \texttt{accuracy\_score} hingga \texttt{confusion\_matrix}, serta model-model algoritma \texttt{LogisticRegression}, \texttt{KNeighborsClassifier}, \texttt{DecisionTreeClassifier}, \texttt{RandomForestClassifier}, dan \texttt{AdaBoostClassifier}) dan modul \texttt{imblearn}.
\end{itemize}

\subsection{Load Dataset}
\begin{figure}[H]
    \centering
    \includegraphics[width=0.7\textwidth]{2.png}
    \caption{Source Code dan Output Load Dataset}
    \label{fig:2}
\end{figure}
Penjelasan Gambar :
\begin{itemize}
    \item \textbf{Input:} Menggunakan library \texttt{gdown} untuk mengunduh file dengan \texttt{file\_id} \textit{'1PnYHfDwq9NjLFLbvbuO9RZQ5mAga838m'} ke dalam file \texttt{healthcare-dataset-stroke-data.csv}. Selanjutnya file tersebut dibaca ke DataFrame \texttt{df}.
    \item \textbf{Output:} Menampilkan baris log \textit{Downloading...} yang menandakan unduhan sebesar 317kB telah sukses (100\%).
\end{itemize}

\subsection{Exploratory Data Analysis (EDA)}

\subsubsection{A. df.tail() dan B. df.head()}
\begin{figure}[H]
    \centering
    \includegraphics[width=0.9\textwidth]{3.png}
    \caption{Output fungsi tail dan head}
    \label{fig:3}
\end{figure}
Penjelasan Gambar :
\begin{itemize}
    \item \textbf{Input:} Menjalankan \texttt{df.tail()} dan \texttt{df.head()} pada dataset.
    \item \textbf{Output:} Menampilkan tabel yang berisi 5 baris terakhir (indeks 5105 hingga 5109) dan 5 baris pertama (indeks 0 hingga 4) dari dataset. Terlihat fitur-fitur yang ada, seperti \texttt{id}, \texttt{gender}, \texttt{age}, \texttt{bmi}, dll, serta nilai kosong (NaN) yang nampak pada kolom \texttt{bmi} di baris indeks 1.
\end{itemize}

\subsubsection{C. Deskripsi Statistik}
\begin{figure}[H]
    \centering
    \includegraphics[width=0.8\textwidth]{4.png}
    \caption{Output df.describe()}
    \label{fig:4}
\end{figure}
Penjelasan Gambar :
\begin{itemize}
    \item \textbf{Input:} Memanggil \texttt{df.describe()} untuk merangkum fitur numerik.
    \item \textbf{Output:} Menampilkan jumlah (count), rata-rata, standar deviasi, min, hingga maks. Jumlah baris data adalah 5110, namun jumlah pada kolom \texttt{bmi} hanya 4909, menandakan kolom tersebut tidak lengkap secara data. Rata-rata umur (\texttt{age}) adalah 43.22.
\end{itemize}

\subsubsection{D. Presentase Target Kelas}
\begin{figure}[H]
    \centering
    \includegraphics[width=0.6\textwidth]{5.png}
    \caption{Source Code dan Output Pie Chart}
    \label{fig:5}
\end{figure}
Penjelasan Gambar :
\begin{itemize}
    \item \textbf{Input:} Melakukan \texttt{value\_counts()} pada target \texttt{stroke} lalu memplotnya sebagai diagram lingkaran (\textit{pie chart}). Terdapat kesalahan penulisan kurung tutup pada \texttt{plt.show}.
    \item \textbf{Output:} Grafik pie chart yang mendemonstrasikan bahwa data sangat tidak seimbang: 95.13\% bernilai 0 (tidak stroke) berbanding 4.87\% yang bernilai 1.
\end{itemize}

\subsubsection{E. Histogram}
\begin{figure}[H]
    \centering
    \includegraphics[width=0.7\textwidth]{6.png}
    \caption{Output Histogram Numerik}
    \label{fig:6}
\end{figure}
Penjelasan Gambar :
\begin{itemize}
    \item \textbf{Input:} Menjalankan \texttt{df.hist(figsize=(10,10))} dan \texttt{plt.show()}.
    \item \textbf{Output:} Diagram histogram untuk variabel numerik (\texttt{id}, \texttt{age}, \texttt{hypertension}, \texttt{heart\_disease}, \texttt{avg\_glucose\_level}, \texttt{bmi}, \texttt{stroke}). Distribusi kolom \texttt{avg\_glucose\_level} terlihat melandai tajam ke sisi kanan.
\end{itemize}

\subsubsection{F. Mengecek Missing Value}
\begin{figure}[H]
    \centering
    \includegraphics[width=0.4\textwidth]{7.png}
    \caption{Output Pengecekan df.isnull().sum()}
    \label{fig:7}
\end{figure}
Penjelasan Gambar :
\begin{itemize}
    \item \textbf{Input:} Memanggil fungsi \texttt{df.isnull().sum()}.
    \item \textbf{Output:} Terbukti secara spesifik bahwa hanya atribut \texttt{bmi} yang memiliki nilai kosong, yakni sebanyak 201 data.
\end{itemize}

\subsection{Data Preprocessing}
\begin{figure}[H]
    \centering
    \includegraphics[width=0.8\textwidth]{8.png}
    \caption{Source Code Data Preprocessing}
    \label{fig:8}
\end{figure}
Penjelasan Gambar :
\begin{itemize}
    \item \textbf{Input:} Kode untuk (A) Memisahkan variabel target \texttt{df\_y} dan prediktor \texttt{df\_X} dengan membuang kolom \texttt{id} dan \texttt{stroke}. (B) Melakukan \texttt{LabelEncoder} pada variabel target, serta melakukan imputasi nilai kosong pada kolom \texttt{bmi} menggunakan median dari kolom tersebut. (C) Mengecek kolom bertipe 'object' dan 'bool' (muncul \textit{output}: \texttt{gender}, \texttt{ever\_married}, \texttt{work\_type}, \texttt{Residence\_type}, \texttt{smoking\_status}) dan mengenkodenya satu per satu. (D) Memecah data latih-uji (\texttt{train\_test\_split}) 70:30 lalu melakukan penyekalaan (\texttt{StandardScaler}).
\end{itemize}

\subsection{Pemeriksaan}
\begin{figure}[H]
    \centering
    \includegraphics[width=0.6\textwidth]{9.png}
    \caption{Output Tinjauan Variabel Preprocessing}
    \label{fig:9}
\end{figure}
Penjelasan Gambar :
\begin{itemize}
    \item \textbf{Input:} Mengeksekusi nama variabel \texttt{X}, \texttt{y}, \texttt{X\_train}, dan \texttt{X\_test} pada sel.
    \item \textbf{Output:} Matriks numerik (numpy arrays) hasil transformasi data sebelumnya, memperlihatkan nilai desimal fitur yang telah melewati standardisasi.
\end{itemize}

\subsection{Modelling}

\subsubsection{A. Logistic Regression}
\begin{figure}[H]
    \centering
    \includegraphics[width=0.8\textwidth]{10.png}
    \caption{Source Code Logistic Regression}
    \label{fig:10}
\end{figure}
\begin{figure}[H]
    \centering
    \begin{minipage}{0.45\textwidth}
        \centering
        \includegraphics[width=\textwidth]{11.png}
        \caption{Output Evaluasi Log-Reg}
        \label{fig:11}
    \end{minipage}\hfill
    \begin{minipage}{0.45\textwidth}
        \centering
        \includegraphics[width=\textwidth]{12.png}
        \caption{Output Koefisien dan Intercept}
        \label{fig:12}
    \end{minipage}
\end{figure}
Penjelasan Gambar :
\begin{itemize}
    \item \textbf{Input:} Baris kode pemodelan \texttt{LogisticRegression()}, fitting dengan data \texttt{x\_train}, pemanggilan prediksi, serta kode visualisasi CM (\textit{Confusion Matrix}) menggunakan \texttt{ConfusionMatrixDisplay} (Gambar 10). Terdapat juga kode memanggil \texttt{model.coef\_} dan \texttt{model.intercept\_}.
    \item \textbf{Output:} Visualisasi dan teks hasil evaluasi di Gambar 11 menunjukkan \textit{Accuracy} 0.9419. Dari susunan matriks: 1444 True Negatives, 0 False Positives, 89 False Negatives, dan 0 True Positives. Gambar 12 menampilkan matriks koefisien berbobot array (e.g. -0.0187) dan intercept bernilai -4.045.
\end{itemize}

\subsubsection{B. K-Nearest Neighbors}
\begin{figure}[H]
    \centering
    \includegraphics[width=0.8\textwidth]{13.png}
    \caption{Source Code K-Nearest Neighbors}
    \label{fig:13}
\end{figure}
\begin{figure}[H]
    \centering
    \includegraphics[width=0.5\textwidth]{14.png}
    \caption{Output Evaluasi K-Nearest Neighbors}
    \label{fig:14}
\end{figure}
Penjelasan Gambar :
\begin{itemize}
    \item \textbf{Input:} Membentuk model \texttt{KNeighborsClassifier(n\_neighbors=10)}.
    \item \textbf{Output:} Evaluasi kembali menghasilkan matriks CM [1444, 0], [89, 0] dengan \textit{Accuracy} persis di 0.9419, mengindikasikan ketidakmampuan penuh mendeteksi kelas positif.
\end{itemize}

\subsubsection{C. Decision Tree}
\begin{figure}[H]
    \centering
    \includegraphics[width=0.8\textwidth]{15.png}
    \caption{Source Code Decision Tree}
    \label{fig:15}
\end{figure}
\begin{figure}[H]
    \centering
    \includegraphics[width=0.5\textwidth]{16.png}
    \caption{Output Evaluasi Decision Tree}
    \label{fig:16}
\end{figure}
Penjelasan Gambar :
\begin{itemize}
    \item \textbf{Input:} Inisialisasi model \texttt{DecisionTreeClassifier(criterion="entropy")}.
    \item \textbf{Output:} Hasil CM menunjukan perubahan menjadi: 1374 TN, 70 FP, 77 FN, dan 12 TP. Akurasi turun menjadi 0.9041, namun True Positives berhasil dideteksi meski kecil (12).
\end{itemize}

\subsubsection{D. Random Forest}
\begin{figure}[H]
    \centering
    \includegraphics[width=0.8\textwidth]{17.png}
    \caption{Source Code Random Forest}
    \label{fig:17}
\end{figure}
\begin{figure}[H]
    \centering
    \includegraphics[width=0.5\textwidth]{18.png}
    \caption{Output Evaluasi Random Forest}
    \label{fig:18}
\end{figure}
Penjelasan Gambar :
\begin{itemize}
    \item \textbf{Input:} Inisialisasi \texttt{RandomForestClassifier(random\_state=42)}.
    \item \textbf{Output:} Matriks CM menunjukkan: 1442 TN, 2 FP, 89 FN, dan 0 TP. Akurasi berada di 0.9406 dengan mayoritas tetap terprediksi sebagai 0.
\end{itemize}

\subsubsection{E. AdaBoost}
\begin{figure}[H]
    \centering
    \includegraphics[width=0.8\textwidth]{19.png}
    \caption{Source Code AdaBoost}
    \label{fig:19}
\end{figure}
\begin{figure}[H]
    \centering
    \includegraphics[width=0.5\textwidth]{20.png}
    \caption{Output Evaluasi AdaBoost}
    \label{fig:20}
\end{figure}
Penjelasan Gambar :
\begin{itemize}
    \item \textbf{Input:} Inisialisasi \texttt{AdaBoostClassifier(random\_state=42)}.
    \item \textbf{Output:} Hasil evaluasi CM sama persis dengan Logistic Regression, yaitu [1444, 0] dan [89, 0], berakurasi 0.9419.
\end{itemize}


% ==========================================
% TUGAS PRAKTIKUM
% ==========================================
\clearpage
\section{Tugas Praktikum}

\subsection{Impor Pustaka}
\begin{figure}[H]
    \centering
    \includegraphics[width=0.8\textwidth]{T1.png}
    \caption{Source Code Impor Pustaka}
    \label{fig:t1}
\end{figure}
Penjelasan Gambar :
\begin{itemize}
    \item \textbf{Input:} Memasukkan modul \texttt{pandas}, \texttt{numpy}, \texttt{matplotlib.pyplot}, \texttt{sklearn}, dan \texttt{imblearn} ke dalam sistem sebelum memulai proses analisis dan klasifikasi.
\end{itemize}

\subsection{Load Dataset dan Melihat Baris Data}
\begin{figure}[H]
    \centering
    \includegraphics[width=0.8\textwidth]{T2.png}
    \caption{Output Load Dataset, df.head() dan df.tail()}
    \label{fig:t2}
\end{figure}
Penjelasan Gambar :
\begin{itemize}
    \item \textbf{Input:} Mengunduh dan membaca \texttt{loan-data.csv}, dilanjutkan dengan eksekusi fungsi \texttt{df.head()} dan \texttt{df.tail()} untuk meninjau baris teratas serta baris paling bawah.
    \item \textbf{Output:} Menampilkan isi struktur tabel dari \texttt{Loan\_ID} hingga \texttt{Loan\_Status}.
\end{itemize}

\subsection{Deskripsi Statistik}
\begin{figure}[H]
    \centering
    \includegraphics[width=0.9\textwidth]{T3.png}
    \caption{Output df.describe()}
    \label{fig:t3}
\end{figure}
Penjelasan Gambar :
\begin{itemize}
    \item \textbf{Input:} Menjalankan perintah \texttt{df.describe()}.
    \item \textbf{Output:} Rekapitulasi distribusi data numerik yang memberikan ringkasan \textit{mean}, standar deviasi, hingga angka minimum dan maksimum.
\end{itemize}

\subsection{Presentase Target dan Histogram}
\begin{figure}[H]
    \centering
    \includegraphics[width=0.7\textwidth]{T4.png}
    \caption{Source Code Visualisasi Target dan Histogram}
    \label{fig:t4}
\end{figure}
\begin{figure}[H]
    \centering
    \begin{minipage}{0.45\textwidth}
        \centering
        \includegraphics[width=\textwidth]{T5.png}
        \caption{Output Hasil Pie Chart}
        \label{fig:t5}
    \end{minipage}\hfill
    \begin{minipage}{0.45\textwidth}
        \centering
        \includegraphics[width=\textwidth]{T6.png}
        \caption{Output Hasil Histogram}
        \label{fig:t6}
    \end{minipage}
\end{figure}
Penjelasan Gambar :
\begin{itemize}
    \item \textbf{Input:} Blok kode (T4) untuk mengonversi jumlah kelas pada \texttt{Loan\_Status} menjadi grafik bundar dan pembuatan diagram histogram untuk variabel numerik.
    \item \textbf{Output:} (T5) Menampilkan \textit{pie chart} di mana kelas target mayoritas adalah Y (disetujui). (T6) Plot grid histogram memperlihatkan persebaran masing-masing data fitur numerik.
\end{itemize}

\subsection{Pengecekan Missing Value}
\begin{figure}[H]
    \centering
    \includegraphics[width=0.4\textwidth]{T7.png}
    \caption{Output df.isnull().sum()}
    \label{fig:t7}
\end{figure}
Penjelasan Gambar :
\begin{itemize}
    \item \textbf{Input:} Perintah \texttt{df.isnull().sum()}.
    \item \textbf{Output:} Daftar jumlah data yang hilang di masing-masing kolom dataset, membantu penentuan jenis \textit{imputation} pada langkah pemrosesan.
\end{itemize}

\subsection{Data Preprocessing}
\begin{figure}[H]
    \centering
    \includegraphics[width=0.8\textwidth]{T8.png}
    \caption{Source Code Data Preprocessing}
    \label{fig:t8}
\end{figure}
Penjelasan Gambar :
\begin{itemize}
    \item \textbf{Input:} Keseluruhan tahapan pembersihan data: mengisi \textit{missing values}, menghapus atribut id, \textit{Label Encoding} untuk data string/teks, pemisahan data latih/uji (70:30), serta penerapkan \texttt{StandardScaler} untuk menormalisasi nilai numerik.
\end{itemize}

\subsection{Pemeriksaan Variabel}
\begin{figure}[H]
    \centering
    \begin{minipage}{0.45\textwidth}
        \centering
        \includegraphics[width=\textwidth]{T9.png}
        \caption{Output Pengecekan Variabel 1}
        \label{fig:t9}
    \end{minipage}\hfill
    \begin{minipage}{0.45\textwidth}
        \centering
        \includegraphics[width=\textwidth]{T10.png}
        \caption{Output Pengecekan Variabel 2}
        \label{fig:t10}
    \end{minipage}
\end{figure}
Penjelasan Gambar :
\begin{itemize}
    \item \textbf{Output:} (T9 \& T10) Array numpy dari fitur (X) maupun target (y) termasuk data uji serta data latih, memastikan semuanya telah bernilai numerik sesuai proses transformasi sebelumnya.
\end{itemize}

\subsection{Modelling}
\begin{figure}[H]
    \centering
    \includegraphics[width=0.8\textwidth]{T11.png}
    \caption{Source Code Definisi Model}
    \label{fig:t11}
\end{figure}
Penjelasan Gambar :
\begin{itemize}
    \item \textbf{Input:} Pendefinisian blok fungsi-fungsi yang menginisialisasi serta melatih (fitting) lima model algoritma utama (Logistic Regression, KNN, Decision Tree, Random Forest, AdaBoost).
\end{itemize}

\subsection{Menjalankan Model}
\begin{figure}[H]
    \centering
    \includegraphics[width=0.8\textwidth]{T12.png}
    \caption{Source Code Run Model}
    \label{fig:t12}
\end{figure}
Penjelasan Gambar :
\begin{itemize}
    \item \textbf{Input:} Memanggil semua fungsi pemodelan tersebut sekaligus agar tereksekusi, dan menyimpan nilai rekap hasil evaluasi tiap model ke dalam sebuah \textit{list} yang akan diproses menjadi data tabular.
\end{itemize}

\subsection{Hasil Confusion Matrix per Model}
\begin{figure}[H]
    \centering
    \begin{minipage}{0.45\textwidth}
        \centering
        \includegraphics[width=\textwidth]{T13.png}
        \caption{CM Logistic Regression}
        \label{fig:t13}
    \end{minipage}\hfill
    \begin{minipage}{0.45\textwidth}
        \centering
        \includegraphics[width=\textwidth]{T14.png}
        \caption{CM K-Nearest Neighbors}
        \label{fig:t14}
    \end{minipage}
\end{figure}

\begin{figure}[H]
    \centering
    \begin{minipage}{0.32\textwidth}
        \centering
        \includegraphics[width=\textwidth]{T15.png}
        \caption{CM Decision Tree}
        \label{fig:t15}
    \end{minipage}\hfill
    \begin{minipage}{0.32\textwidth}
        \centering
        \includegraphics[width=\textwidth]{T16.png}
        \caption{CM Random Forest}
        \label{fig:t16}
    \end{minipage}\hfill
    \begin{minipage}{0.32\textwidth}
        \centering
        \includegraphics[width=\textwidth]{T17.png}
        \caption{CM AdaBoost}
        \label{fig:t17}
    \end{minipage}
\end{figure}
Penjelasan Gambar :
\begin{itemize}
    \item \textbf{Output:} Matriks plot \textit{Confusion Matrix} hasil tebakan untuk masing-masing algoritma klasifikasi (T13: LogReg, T14: KNN, T15: Decision Tree, T16: Random Forest, T17: AdaBoost) dalam mendeteksi True Positive, True Negative, False Positive, dan False Negative pada data uji.
\end{itemize}

\subsection{Perbandingan Model}
\begin{figure}[H]
    \centering
    \includegraphics[width=0.9\textwidth]{T18.png}
    \caption{Output Tabel Perbandingan}
    \label{fig:t18}
\end{figure}
Penjelasan Gambar :
\begin{itemize}
    \item \textbf{Output:} Tabel rekap hasil perbandingan akhir (F1 Score, Accuracy, Precision, Recall). Menunjukkan evaluasi menyeluruh untuk menentukan seberapa optimal setiap model dalam menangani masalah persetujuan pinjaman ini.
\end{itemize}

\chapter{Kesimpulan}
\label{chap:kesimpulan}
Berdasarkan praktikum mengenai klasifikasi pada \textit{machine learning} yang telah dilaksanakan—dimulai dari tahap persiapan data hingga pemodelan komparatif berbagai algoritma—diperoleh beberapa simpulan penting sebagai berikut:

\begin{enumerate}
    \item \textbf{Pentingnya Data Preprocessing sebelum Pemodelan:} Proses \textit{Exploratory Data Analysis} (EDA) dan \textit{data preprocessing} (seperti penanganan \textit{missing values}, \textit{label encoding} untuk fitur kategorial, serta \textit{standard scaling}) terbukti menjadi tahapan esensial. Data yang bersih dan berada dalam skala yang seragam sangat memengaruhi keberhasilan komputasi jarak maupun konvergensi algoritma klasifikasi.
    
    \item \textbf{Performa Algoritma Klasifikasi yang Beragam:} Berbagai model klasifikasi (Logistic Regression, K-Nearest Neighbors, Decision Tree, Random Forest, dan AdaBoost) berhasil diimplementasikan. Hasil eksperimen pada klasifikasi kelayakan pinjaman menunjukkan bahwa algoritma \textit{ensemble}, khususnya \textbf{Random Forest}, memberikan performa yang sedikit lebih unggul (mencapai \textit{Accuracy} 0.8261 dan \textit{F1-Score} makro 0.7515) karena kemampuannya meminimalisasi varians dan risiko \textit{overfitting} dibanding model pohon tunggal (\textit{Decision Tree}).
    
    \item \textbf{Keterbatasan Akurasi pada Data Tidak Seimbang (Imbalanced Data):} Pada kasus prediksi dataset yang sangat tidak seimbang (seperti pada data *Stroke* dengan mayoritas kelas negatif $95.13\%$), metrik \textit{Accuracy} tinggi ($>94\%$) bersifat mengecoh (\textit{misleading}). Model seringkali memprediksi mayoritas kelas negatif namun sangat buruk dalam mendeteksi kelas minoritas (positif), sehingga evaluasi lanjutan sangat diperlukan.
    
    \item \textbf{Peran Vital Confusion Matrix dan Metrik Lanjutan:} Penggunaan \textit{Confusion Matrix} memberikan transparansi nyata terhadap distribusi prediksi benar (TP, TN) maupun galat (FP, FN). Hal ini membuktikan bahwa pengukuran kualitas model tidak cukup hanya dari \textit{Accuracy}, melainkan wajib mempertimbangkan \textit{Precision}, \textit{Recall}, serta \textit{F1-Score} untuk mendapatkan keseimbangan prediksi yang komprehensif.
    
    \item \textbf{Ekosistem Scikit-Learn yang Modular:} Integrasi pustaka \texttt{scikit-learn} di Python terbukti mempermudah keseluruhan alur kerja (pipeline) klasifikasi. Mulai dari prapemrosesan data (\texttt{StandardScaler}, \texttt{LabelEncoder}), pemodelan berbagai metode pengklasifikasi, hingga kalkulasi evaluasi metrik dapat dilakukan secara instan, modular, dan efisien.
\end{enumerate}

\chapter{Daftar Pustaka}
\begin{enumerate}
    \item \label{modulPraktikum} Dr. Eng. Ir. Ganjar Alfian, S.T., M. Eng., dan Tim Asisten. \textit{Modul Praktikum Penambangan Data: Exploratory Data Analysis (EDA) dan Data Preprocessing}. Yogyakarta: Departemen Teknik Elektro dan Informatika, Sekolah Vokasi, Universitas Gadjah Mada, 2026.
    
    \item \label{tukey} J. W. Tukey. \textit{Exploratory Data Analysis}. Reading, MA: Addison-Wesley, 1977.
    
    \item \label{mckinney} W. McKinney. \textit{Python for Data Analysis: Data Wrangling with Pandas, NumPy, and Jupyter}, 3rd ed. Sebastopol, CA: O'Reilly Media, 2022.
    
    \item \label{geron} A. Géron. \textit{Hands-On Machine Learning with Scikit-Learn, Keras, and TensorFlow}, 3rd ed. Sebastopol, CA: O'Reilly Media, 2022.
    
    \item \label{bishop} C. M. Bishop. \textit{Pattern Recognition and Machine Learning}. New York, NY: Springer, 2006.
    
    \item \label{kaggleStroke} fedesoriano. \textit{Stroke Prediction Dataset}. Kaggle, 2021. [Online]. Tersedia: \url{https://www.kaggle.com/datasets/fedesoriano/stroke-prediction-dataset}.
    
    \item \label{kaggleLoan} Bhavik Jikadara. \textit{Loan Status Prediction}. Kaggle, 2024. [Online]. Tersedia: \url{https://www.kaggle.com/datasets/bhavikjikadara/loan-status-prediction}.
    
    \item \label{colabLaporan} H. C. A. Paliling. \textit{S5-PPD-542159 (Source Code Laporan Praktikum)}. Google Colaboratory, 2026. [Online]. Tersedia: \url{https://colab.research.google.com/drive/1bvBJJTwzyzKHsH95cHhvi076hqc8DzNy?usp=sharing}.
\end{enumerate}
